% content of the chapter relativity
\chapter{Relativity}
\label{relativity}

Prior to 1905, physics was dominated by Newtonian mechanics and Maxwell’s 
equations of electromagnetism. However, a conflict existed between 
the two. Newtonian physics allowed for infinite speeds, while Maxwell’s 
equations suggested that light travels at a fixed speed, $c$.
The "luminiferous aether" --- a hypothetical medium thought to carry 
light wave --- was proposed to resolve this, but the famous Michelson-Morley
experiment failed to detect it. Albert Einstein resolved this paradox 
not by fixing the old theories, but by reinventing the definitions of space and time.

Special Relativity  is built upon two deceptively simple postulates:
\begin{description}
    \item [The Principle of Relativity] 
       The laws of physics are the same in all inertial frames of reference.
   \item[The Constancy of the Speed of Light] 
       The speed of light in a vacuum ($c \approx 3.00 \times 10^8$ m/s) is constant for all observers, regardless of the motion of the light source or the observer.
\end{description}

If the speed of light must remain constant, then space and time must be flexible to accommodate it.
\begin{description}
    \item[Time Dilation] Time is not absolute. A moving clock ticks more slowly than a stationary clock. As an object approaches the speed of light, time slows down significantly relative to a stationary observer.
 \item[Length Contraction] Objects moving at relativistic speeds appear shorter along the direction of motion to a stationary observer.
 \item[Mass-Energy Equivalence] 
    $E = mc^2$
    This implies that a small amount of mass can be converted into a tremendous amount of energy, the principle behind nuclear power and stars.
\end{description}
General relativity adds time dilation due to gravity to this list.

There is solid  experimental evidence in favor of relativity, both 
from specially designed physical experiments and the facts of everyday life.
The Global Positioning System (GPS) is a daily application of relativity.
 Satellites move fast ($\approx 14,000$ km/h), causing their clocks to tick {\em slower} by 7 microseconds per day according to Special Relativity.
 Satellites are far from Earth's mass (weaker gravity), causing their clocks to tick {\em faster} by 45 microseconds per day according to General Relativity.
Without correcting for this net difference of 38 microseconds daily, GPS accuracy would drift by kilometers every day, rendering the system useless.

The conventional wisdom that classical mechanics is recovered from relativity 
when $v/c \ll 1$ is deeply entrenched in the physics literature. Textbooks 
frequently state that Newtonian formulas emerge by taking the first terms in 
a Taylor expansion of relativistic expressions, or by taking the limit $c \to \infty$.
While such approximations correctly reproduce certain numerical results 
(for one particle), they do not recover the {\em classical mechanics}.
The absence of absolute simultaneity, the lack of global time, incompatibility 
of instantaneous constraints, and the necessity of fields for 
interactions signal a profound structural divergence.

Newtonian mechanics requires a global time parameter shared by all 
degrees of freedom. This time orders events, defines simultaneity, 
and underlies the Hamiltonian formulation. Phase space trajectories 
evolve with respect to a universal time coordinate. The Newtonian $N$-particle 
configuration space $\mathbb{R}^{3N}$ is defined on simultaneous positions of all particles.

Rigid bodies in classical mechanics require instantaneous constraint 
enforcement and a notion of shape preserved by all inertial observers. 
These presuppose global simultaneity. In relativity, Born rigidity 
is possible only under highly restricted motions and cannot be maintained 
for general accelerations. Constraint forces propagate at finite speed, 
so multi-point instantaneous constraints are prohibited.

Constraints defining rigid links, rods, joints, or holonomic relations 
are not relativistically admissible, except as approximations mediated 
through elastic or electromagnetic fields with finite propagation speed. 
The configuration spaces associated with classical mechanical systems 
(e.g., rigid rotors, articulated bodies) have no relativistic analogue.

The famous no-interaction theorem (\cite{cjs-no-int},\cite{leutwyler}) 
demonstrates that a system of finitely many relativistic particles 
cannot interact through potentials if one demands Poincar'e invariance 
and a single time parameter. Nontrivial interactions require fields. 
Thus classical particle mechanics with inter particle forces has no relativistic extension.

This result explains the structural collapse of the Newtonian 
$N$-body problem in relativity. The notion of an instantaneous 
force depending on simultaneous positions is meaningless. 
Relativity allows only local interactions mediated by fields or 
direct interactions (collisions).

One might hope that classical continuum mechanics survives relativity 
even if rigid bodies do not. However, relativistic continuum mechanics 
is fundamentally different from its Newtonian counterpart. 
A relativistic continuum is defined by a stress--energy tensor $T^{ij}$ 
satisfying local conservation laws. There is no concept of a 
perfectly rigid body, and the equations of state must respect causality.

In Newtonian continuum mechanics, pressure and stress can adjust 
instantaneously across the medium. In relativity, signals propagate 
at finite speed, invalidating classical constitutive relations. 
The Newtonian limit of relativistic continua yields compressible, 
not rigid, media. The structural assumptions are thus incompatible.

Given the collapse of constraints, multi-particle structure, 
and continuum rigidity, the only part of classical mechanics 
that survives intact in relativity is the motion of a single 
particle under an external electromagnetic field. 
The Lorentz force law is compatible with Minkowski geometry and 
requires no global time. Maxwell's equations provide the necessary field mediation.
Everything else---interacting particles, rigid bodies, 
constrained systems, Newtonian continua---is not contained within relativity.

Classical mechanics is not a limit case of relativity. 
These are two different models of reality with one particle
problems as common domain and they should cooperate for
physical problems. It is not uncommon that only part of
the system is relativistic, e.g. fast electrons may
interact with slow ions. Besides that for most physical
problems {\em there is} natural frame, the most convenient
for analysis.

\section{Special Relativity}
\label{special-relativity}

\subsection{Relativistic particle}
\label {relativistic-particle}

 Lagrangian function for relativistic
particle in electromagnetic field can be written as
\cite{Landau-Lifshitz-2},\cite{Goldstein}:
\begin{equation}
   \label {lagr-3d}
   L=-m c^2 \sqrt{1-\frac{v^2}{c^2}} - e \phi + \frac{e}{c} \vec A \vec v
\end{equation}
So canonical moment $ \vec P = \partial L/\partial \vec v $ is
\begin{equation}
   \label {mon-3d}
    \vec P=\frac{m \vec v}{\sqrt{1-\frac{v^2}{c^2}}} +\frac{e}{c} \vec A 
\end{equation}
and Hamiltonian is
\begin{equation}
   \label {ham-3d}
    H=\sqrt{m^2 c^4+c^2 (\vec P -\frac{e}{c} \vec A)^2} +e \phi
\end{equation}
and equation of motion are
\begin{equation}
   \label {eqm-3d}
   \frac{d \vec P}{d t}=e \vec E + \frac{e}{c} \vec v \times \vec B
\end{equation}

\subsection{Minkowski spacetime}
\label {minkowski-spacetime}
Newtonian mechanics is inseparable from Galilean spacetime,
which possesses
absolute time and Euclidean space. The geometry consists of a 
three-dimensional spatial manifold evolving in a universal time $t$.
Simultaneity is observer-independent. The symmetry group is 
the Galilei group, comprising translations,rotations and Galilei boost
(see Chapter \ref{class-mech}):
\begin{equation}
   \label {gal-boost}
   \begin{split}
       t' &= t \\
       \vec x' &= \vec x -\vec v t
   \end{split}
\end{equation}

In contrast, special relativity is formulated in Minkowski spacetime (see \ref{me-cov})
 and the symmetry group is the Poincare group, consisting of
 inhomogeneous linear transformations preserving metric \eqref{minkg}, which also
 contains translations and rotations, but Lorenz boost
 \begin{equation}
    \label {lor-boost}
    \begin{split}
        t' &= \gamma (t-\frac{\vec v \vec x}{c^2}) \\
        \vec x' &= \vec x + \frac{(\gamma -1)}{v^2} (\vec v \cdot \vec x) \vec v-\gamma \vec v t  \\
        \gamma &= \frac{1}{\sqrt{1-\frac{v^2}{c^2}}}
    \end{split}
 \end{equation}
 replaces Galilean one.

Points in that manifold are called events. Taking arbitrary 
event as the origin one can see that
The separation between events (square of invariant interval) 
\begin{equation}
   \label {interval}
    s^2=c^2 t^2-|\vec x|^2
\end{equation}
may be timelike ($s>0$), lightlike ($s=0$) or spacelike ($s<0$).
So when the origin in Minkowski space is fixed, it is divided into
five regions: future ($ s>0 $ $ t>0 $), cone ($ s>0 $ $ t<0 $),
future light cone ($ s=0 $ $ t>0 $), past light cone ($ s=0 $ $ t<0 $)
and elsewhere ($ s<0 $).

The limit $c \to \infty$ is often said to make \eqref{gal-boost} 
from \eqref{lor-boost} 
and collapse the light cones and elsewhere 
into a spatial hyperplane $ t=0 $ --- now. However this limit
does not exists in any mathematical sense. 

There are three types of vectors in Minkowski space,
a vector $v^i$ is timelike if $g(v,v)= g_{ij} v^i v^j=v_i v^i>0$,
lightlike if $v_i v^i=0$ and spacelike if $v_i v^i<0$
Note that for pair of orthogonal vectors only two possibilities exist:
one of them is timelike another is spacelike, or both are lightlike.


It is often said that relativity makes time and space coordinate equal. 
This is plain wrong,
time is not at the least equal to space position, beside difference in sign there is fundamental contrast: we can always stop the particle so
that is coordinates $ \vec x=0 $ but we can not stop time flow. The only documented case \footnote {
Joshua 10:12-14
Then spake Joshua to the LORD in the day when the LORD delivered 
up the Amorites before the children of Israel, and he said in the 
sight of Israel, Sun, stand thou still upon Gibeon; and thou, 
Moon, in the valley of Ajalon.
And the sun stood still, and the moon stayed, until the people 
had avenged themselves upon their enemies. Is not this written 
in the book of Jasher? So the sun stood still in the midst of heaven,
and hasted not to go down about a whole day.} 
raises doubts about its authenticity.

Any  particle of positive mass has its worldline $x^i(s)$ a curve in $M$ 
with timelike tangent vector. Curves with lightlike tangent vector 
correspond to zero mass objects --- light rays, for example.  

Parameter $ s $ may be arbitrary, for timelike curves we choose it to be the
length of the curve.
\begin{equation}
   \label {int4}
   ds^2=c^2 dt^2-(dx^2+dy^2+dz^2)
\end{equation}
$ d \tau = ds/c $ is called proper time.
We define four-velocity $u^i = \frac{dx^i}{ds}$.
\begin{equation}
   \label {vel4}
    u^i=(\gamma,\gamma \vec v/c)
\end{equation}
Note that 
\begin{equation}
   \label {4vnorm}
u_i u^i=1
\end{equation} 
Four-acceleration is defined as 
\begin{equation}
   \label {4acc}
 w^i = \frac{du^i}{ds}  
\end{equation}
From \eqref{4vnorm} we see that $ u_i w^i=0 $.

Four-momentum is defined as $p^i = m c u^i$,hence $ p^i p_i = m^2 c^2$.
Dynamics is encoded by the equation $\frac{dp^i}{d\tau} = f^i$, note that $ u_i f^i = 0$ (mass-shell preservation).

We can formally expand Principle of Least Action on relativity. Luck of global time 
deprives state of the system in fixed time from any physical sense so limits of integration
become undefined. Nevertheless it remains good way of defining equation of motion.

Covariant action my be written as  
\begin{equation}
   \label {act-1}
    S_1=-m c \int (\sqrt{g_{i j} d x^i d x^j}  - \frac{e}{c} A_i d x^i)
\end{equation}
it is evidently independent of parametrization of the worldline. 
Let $ \dot x^i (\lambda)=d x^i / d \lambda $, $ \lambda $ is 
arbitrary. So Lagrangian is:
\begin{equation}
   \label {lag-1}
    L_1=-m c \sqrt{g_{i j} \dot x^i \dot x^j} -\frac{e}{c} A_i \dot x^i
\end{equation}
momentum 
\begin{equation}
   \label {mom-1}
   P_i = -m c \frac{\dot x_i}{\sqrt{\dot x^l \dot x_l}} -\frac{e}{c} A_i
\end{equation}
and equations of motion are
\begin{equation}
   \label {eqm-1}
    m c \frac{\ddot x_j}{\sqrt{\dot x^l \dot x_l}} (\delta_i^j -
    \frac{\dot x^i \dot x_j}{\dot x^l \dot x_l}) = \frac{e}{c} F_{j i} \dot x^i
\end{equation}
The equation \eqref{eqm-1} does not include longitudinal component of 
4-velocity. If we choose it to satisfy \eqref{4vnorm} (i.e. use natural
parametrisation we get:
\begin{equation}
   \label {eqm-2}
    m c w_j = \frac{e}{c} F_{j i} u^i
\end{equation}
which corresponds to momenta:
\begin{equation}
   \label {mom-2}
   P_i = -m c  u_i -\frac{e}{c} A_i
\end{equation}
and Lagrangian:
\begin{equation}
   \label {lag-2}
    L_2=- \frac{m c}{2} g_{i j} \dot x^i \dot x^j -\frac{e}{c} A_i \dot x^i
\end{equation}
We ignore restriction \eqref{4vnorm} calculating variation of action. 
Though it often said that \eqref{4vnorm} is constraint it is not
actually, it is constant of motion and should be added to \eqref{eqm-2}
as restriction on initial conditions, to provide physically meaningful
solutions. The same situation exactly as with Lorenz gauge for wave
equation \eqref{me-potc}.

\section{General Relativity}
\label{general_relativity}
Special Relativity had a major limitation: it did not account 
for acceleration or gravity. After a decade of struggle, 
Einstein published General Relativity, which redefined 
gravity not as a force, but as the curvature of spacetime.

Newton viewed gravity as a force pulling objects together. 
Einstein viewed it as geometry. Massive objects (like the Sun) 
warp the four-dimensional fabric of spacetime 
(three dimensions of space + one dimension of time).
"Matter tells space how to curve; space tells matter how to move."(John Archibald Wheeler).

Planets do not orbit the Sun because the Sun pulls them; 
they orbit because they are following the straightest possible path 
(a geodesic) along the curved surface of spacetime created by the Sun's mass.
Actually general relativity gives only small corrections even 
at Solar system scale it is probably the most beautiful theory but too
complicated to be considered here. 


